% =====================================================================
%  EnKF Workshop 2026 -- Rosendal, Norway
%  Talk 7: Ensemble Smoothing for Joint State and Parameter Estimation
%          in Turbulent Urban Flows
%  Nikolaj T. Mucke, Femke Vossepoel -- TU Delft  (UrbanAIR project)
%
%  >>> Compile with XeLaTeX (twice for TikZ "remember picture"):
%        xelatex main.tex ; xelatex main.tex
%
%  Structure: ~15 content slides + backup appendix.
%  The talk follows Geir Evensen's invited talk (Talk 6), which
%  introduces the LBM-LES model and the inflow-BC ESMDA setup --
%  this deck deliberately builds on it instead of repeating it.
% =====================================================================
\documentclass[aspectratio=169,11pt]{beamer}

\usepackage{beamerthemeUrbanAIR}
\usepackage{amsmath,amssymb,bm}
\usepackage{booktabs}
\usepackage{array}
\usepackage{textcomp}
\usepackage{colortbl}
\usepackage{animate}      % inline frame animation -- XeLaTeX, no Flash, plays in Adobe Acrobat

% ---- math shortcuts -------------------------------------------------
\newcommand{\R}{\mathbb{R}}
\newcommand{\state}{\bm{x}}
\newcommand{\param}{\bm{\theta}}
\newcommand{\obs}{\bm{d}}
\newcommand{\E}{\mathbb{E}}
\newcommand{\Cov}{\mathbf{C}}

% ---- title metadata -------------------------------------------------
\title{Ensemble Smoothing for Joint State and\\Parameter Estimation in Turbulent Urban Flows}
\subtitle{Challenges of using LES simulators for data assimilation in urban areas}
\author{\textbf{Nikolaj T.\ M\"ucke}, Femke Vossepoel}
\institute{TU Delft \,$\cdot$\, UrbanAIR project}
\setUAevent{EnKF Workshop 2026 \;|\; Rosendal, Norway \;|\; 15 -- 18 June}

\begin{document}

% =====================================================================
%  TITLE
% =====================================================================
\uatitlepage

% =====================================================================
%  OVERVIEW
% =====================================================================
\begin{frame}{Overview}
  \begin{itemize}\setlength\itemsep{11pt}
    \item LES forward models: uDALES, PALM, LBM
    \item The \texttt{pyurbanair} framework
    \item Parameter estimation with ESMDA
    \item State estimation
    \item Neural surrogate model
  \end{itemize}
\end{frame}

% =====================================================================
%  MOTIVATION
% =====================================================================
\begin{frame}{Digital twins for cities}
  \begin{columns}[T]
    \begin{column}{0.58\textwidth}
      \textbf{UrbanAIR} builds digital twins for urban air
      quality and heat (part of \emph{Destination Earth}).
      \smallskip
      \begin{itemize}\setlength\itemsep{6pt}
        \item Street-level forecasts of heat \& pollution hotspots
        \item Decision support: zoning, greening, extreme-heat planning
        \item Uncertainty quantificaiton, data assimilation
      \end{itemize}
    \end{column}
    \begin{column}{0.40\textwidth}
      \uacard{%
        \footnotesize
        \uakeyt{Horizon Europe}\\ \texteuro14.3\,M, 2025--2028\\[4pt]
        \uakeyt{Consortium}\\ 18 partners, 11 countries\\[4pt]
        \uakeyt{Cities}\\ Antwerp, Barcelona;\\ Paris, Bristol,
        Rotterdam}
    \end{column}
  \end{columns}
\end{frame}

% ---- problem recap: builds on Evensen's talk -------------------------
\begin{frame}{Connecting regional model with local model}
  \centering
  \begin{tikzpicture}[font=\scriptsize,>=latex]
    % ===== coarse regional grid (left) =====
    \fill[uaOrange!22] (2.55,0.85) rectangle (3.4,1.70);          % highlighted cell
    \draw[uaCharcoal!45,line width=0.5pt,step=0.85] (0,0) grid (3.4,2.55);
    \draw[uaCharcoal!70,line width=0.9pt] (0,0) rectangle (3.4,2.55);
    \draw[uaOrange,line width=1.1pt] (2.55,0.85) rectangle (3.4,1.70);
    \node[font=\footnotesize\bfseries,text=uaCharcoal] at (1.7,2.95)
      {Regional forecast};
    \node[text=uaGrey,align=center] at (1.7,-0.55)
      {coarse $\Delta x,\ \Delta t$ $\Rightarrow$\\ uncertain inflow prior};

    % ===== fine urban grid (right) =====
    \draw[uaTeal!30,line width=0.3pt,step=0.2] (7.8,0) grid (10.6,2.8);
    \fill[uaTeal] (8.2,0.4) rectangle ++(0.4,0.4);
    \fill[uaTeal] (9.0,0.4) rectangle ++(0.4,0.4);
    \fill[uaTeal] (9.8,0.4) rectangle ++(0.4,0.4);
    \fill[uaTeal] (8.6,1.2) rectangle ++(0.4,0.4);
    \fill[uaTeal] (9.4,1.2) rectangle ++(0.4,0.4);
    \fill[uaTeal] (10.2,1.2) rectangle ++(0.4,0.4);
    \fill[uaTeal] (8.2,2.0) rectangle ++(0.4,0.4);
    \fill[uaTeal] (9.0,2.0) rectangle ++(0.4,0.4);
    \fill[uaTeal] (9.8,2.0) rectangle ++(0.4,0.4);
    \draw[uaTeal,line width=0.9pt] (7.8,0) rectangle (10.6,2.8);
    \node[font=\footnotesize\bfseries,text=uaTeal] at (9.2,3.05)
      {Urban LES + assimilation};
    \node[text=uaGrey,align=center] at (9.2,-0.55)
      {fine $\Delta x,\ \Delta t$ $\Rightarrow$\\ state $\state$ + inflow $\param(t)$};

    % ===== zoom trapezoid =====
    \draw[uaGrey,densely dashed,line width=0.6pt] (3.4,1.70) -- (7.8,2.8);
    \draw[uaGrey,densely dashed,line width=0.6pt] (3.4,0.85) -- (7.8,0.0);

    % ===== downscaling arrow + external-prior chip =====
    \draw[->,line width=1.3pt,uaTeal] (3.6,1.27) -- (7.6,1.4);
    \node[rounded corners=3pt,draw=uaAmber,line width=1pt,fill=uaAmber!16,
          align=center,inner sep=3pt] at (5.6,1.4)
      {\textbf{external prior}\\ $\param_{\text{ext}} \pm \Sigma_{\text{ext}}$};
  \end{tikzpicture}
  \par\medskip
  Correct the \uakey{time-varying inflow} (wind speed \&
  direction) of an urban LES from \uakey{sparse in-domain sensors}.
\end{frame}


% ---- multiscale framing ----------------------------------------------
\begin{frame}{Urban turbulence is multiscale}
  \centering
  \includegraphics[width=0.90\linewidth]{assets/multiscale.png}
\end{frame}


% ---- why hard + what this talk examines ------------------------------
\begin{frame}{Why this is hard}
  \begin{columns}[T]
    \begin{column}{0.32\textwidth}
      \uacard[uaTeal]{\centering
        \uakeyt{Strong nonlinearity}\\[4pt]\footnotesize
        chaotic dynamics; inflow changes cascade through wakes}
    \end{column}
    \begin{column}{0.32\textwidth}
      \uacard[uaOrange]{\centering
        \uakey{Non-Gaussianity}\\[4pt]\footnotesize
        flow--building interactions break Gaussian assumptions}
    \end{column}
    \begin{column}{0.32\textwidth}
      \uacard[uaAmber]{\centering
        {\bfseries\color{uaAmber}High dimension}\\[4pt]\footnotesize
        $\dim(\state)\sim10^{5}$--$10^{7}$,\\ $N_e\sim\mathcal{O}(10^{2})$}
    \end{column}
  \end{columns}
  % \vfill
  % \centering
  % {\large Which \uakey{design choices} make ensemble smoothing work here?}
  % \par\bigskip
  % \begin{columns}[T]
  %   \begin{column}{0.24\textwidth}\centering
  %     \uakeyt{1.\ LES models}\\
  %     {\small three solvers,\\ two paradigms}
  %   \end{column}
  %   \begin{column}{0.24\textwidth}\centering
  %     \uakeyt{2.\ Data assimilation}\\
  %     {\small ES-MDA, windows, localization}
  %   \end{column}
  %   \begin{column}{0.24\textwidth}\centering
  %     \uakeyt{3.\ \texttt{pyurbanair}}\\
  %     {\small one framework\\ + results}
  %   \end{column}
  %   \begin{column}{0.24\textwidth}\centering
  %     \uakeyt{4.\ Neural surrogates}\\
  %     {\small a learned solver\\ + results}
  %   \end{column}
  % \end{columns}
\end{frame}


% \begin{frame}{One benchmark geometry, three CFD solvers}
%   \begin{columns}[T]
%     \begin{column}{0.55\textwidth}
%       % Why run several models on the \uakeyt{same} case?
%       \medskip
%       \begin{itemize}\setlength\itemsep{8pt}
%         \item \uakey{Cross-model twins} 
%         \item \uakey{Fidelity vs.\ cost}
%         \item Clean comparison: same geometry, same DA pipeline
%       \end{itemize}
%     \end{column}
%     \begin{column}{0.45\textwidth}
%       \centering
%       \includegraphics[width=\linewidth]{assets/xie_castro_geometry.png}\\[2pt]
%       {\scriptsize\color{uaGrey} Xie \& Castro (2008): staggered cube
%       array, 25\% plan density -- the standard urban-LES benchmark.}
%     \end{column}
%   \end{columns}
% \end{frame}

% =====================================================================
%  PART 1 -- LES FORWARD MODELS
% =====================================================================
\uasection[Part 1]{LES forward models}

\begin{frame}{Large-Eddy Simulation in a nutshell}
  Resolve the large eddies, model the sub-grid scales. Filtered
  incompressible Navier--Stokes:
  \[
    \frac{\partial \bar{\bm{u}}}{\partial t}
    + (\bar{\bm{u}}\!\cdot\!\nabla)\bar{\bm{u}}
    = -\frac{1}{\rho}\nabla \bar{p}
    + \nu\nabla^{2}\bar{\bm{u}}
    - \nabla\!\cdot\!\bm{\tau}^{\mathrm{sgs}},
    \qquad \nabla\!\cdot\!\bar{\bm{u}} = 0 .
  \]
  \begin{itemize}\setlength\itemsep{3pt}
    \item \uakeyt{sub-grid stress} $\bm{\tau}^{\mathrm{sgs}}$
    \item buildings as no-slip walls
    \item chaotic
  \end{itemize}
\end{frame}


\begin{frame}{Three LES solvers}
  \centering
  \begin{tikzpicture}[font=\footnotesize,
      banner/.style={rounded corners=4pt,align=center,minimum height=1.0cm,
                     inner sep=4pt},
      mdl/.style={rounded corners=4pt,align=center,text=white,
                  minimum height=1.05cm,inner sep=4pt},
      sub/.style={align=center,color=uaGrey,font=\scriptsize}]
    % ---- paradigm divider ----
    \draw[uaOrange,line width=1pt,dash pattern=on 3pt off 2pt] (6.2,-1.55) -- (6.2,2.55);
    % ---- header banners ----
    \node[banner,fill=uaCharcoal!8,text width=5.5cm] at (2.85,2.0)
      {\bfseries\color{uaCharcoal}CONTINUUM\\[1pt]
       {\scriptsize\color{uaGrey}solve Navier--Stokes directly}};
    \node[banner,fill=uaTeal!12,text width=5.5cm] at (9.55,2.0)
      {\bfseries\color{uaTeal}KINETIC\\[1pt]
       {\scriptsize\color{uaGrey}recover N--S as a limit}};
    % ---- model cards ----
    \node[mdl,fill=uaOrange,text width=2.3cm] at (1.6,0.6)
      {\bfseries uDALES};
    \node[mdl,fill=uaAmber,text width=2.3cm] at (4.25,0.6)
      {\bfseries PALM};
    \node[mdl,fill=uaTeal,text width=3.6cm] at (9.55,0.6)
      {\bfseries LBM};
    % ---- one-line characterizations ----
    % \node[sub,text width=2.5cm] at (1.6,-0.95) {geometry-faithful\\urban specialist};
    % \node[sub,text width=2.5cm] at (4.25,-0.95) {high-order,\\feature-rich workhorse};
    % \node[sub,text width=3.6cm] at (9.55,-0.95) {throughput-first,\\built for ensembles};
  \end{tikzpicture}
  \par\medskip
  Different numerics \emph{and} physics.
\end{frame}

% ---- uDALES ----------------------------------------------------------
% \begin{frame}[t]{uDALES}
%   \uacard[uaCharcoal]{\footnotesize
%     \uakeyt{Core.} Finite-difference Boussinesq N--S. \\ 
%     \uakeyt{Geometry:} conservative cut-cell
%     IBM (STL). \\ 
%     \uakeyt{SGS:} Vreman and Smagorinsky.}
  % \vspace{-2pt}
  % \begin{columns}[T]
  %   \begin{column}{0.5\textwidth}
  %     {\bfseries\color{uaTeal}Advantages}\\[-3pt]
  %     {\color{uaTeal}\rule{\linewidth}{0.8pt}}
  %     \begin{itemize}\footnotesize\setlength\itemsep{2pt}
  %       \item conservative, geometry-faithful walls
  %       \item fits real, non-grid-aligned facades
  %       \item non-dissipative, energy-conserving
  %     \end{itemize}
  %   \end{column}
  %   \begin{column}{0.5\textwidth}
  %     {\bfseries\color{uaOrange}Drawbacks}\\[-3pt]
  %     {\color{uaOrange}\rule{\linewidth}{0.8pt}}
  %     \begin{itemize}\footnotesize\setlength\itemsep{2pt}
  %       \item only 2nd-order accurate interior
  %       \item CPU/MPI only -- no production GPU
  %       \item global pressure solve every step (cost)
  %     \end{itemize}
  %   \end{column}
  % \end{columns}
% \end{frame}

% ---- uDALES discretization (visual) ----------------------------------
\begin{frame}{uDALES: staggered C-grid, cut-cell walls}
  \centering
  \begin{tikzpicture}[font=\scriptsize,>=latex]
    % ================= LEFT: staggered Arakawa C-grid =================
    \draw[uaCharcoal!45,line width=0.6pt,step=1.8] (0,0) grid (3.6,3.6);
    \draw[uaCharcoal!70,line width=0.9pt] (0,0) rectangle (3.6,3.6);
    % pressure / scalars at cell centers
    \foreach \x in {0.9,2.7}
      \foreach \y in {0.9,2.7}
        \fill[uaCharcoal] (\x,\y) circle (2.6pt);
    \node[text=uaCharcoal,anchor=west] at (0.98,1.12) {$p$};
    % u on vertical faces (normal to x)
    \foreach \x in {0,1.8,3.6}
      \foreach \y in {0.9,2.7}
        \draw[->,uaOrange,line width=1.1pt] (\x-0.28,\y) -- (\x+0.28,\y);
    \node[text=uaOrange] at (1.8,3.0) {$u$};
    % w on horizontal faces (normal to z)
    \foreach \y in {0,1.8,3.6}
      \foreach \x in {0.9,2.7}
        \draw[->,uaTeal,line width=1.1pt] (\x,\y-0.28) -- (\x,\y+0.28);
    \node[text=uaTeal] at (0.62,1.95) {$w$};
    \node[font=\footnotesize\bfseries,text=uaCharcoal] at (1.8,4.05)
      {Staggered Arakawa C-grid};
    \node[text=uaGrey,align=center,text width=4.6cm] at (1.8,-0.55) {};

    % ================= RIGHT: cut cells ================================
    \begin{scope}[xshift=7.0cm]
      % solid region (right of the slanted facade)
      \fill[uaCharcoal!22] (1.35,3.6) -- (4.5,0.45) -- (4.5,3.6) -- cycle;
      \draw[uaCharcoal!45,line width=0.5pt,step=0.9] (0,0) grid (4.5,3.6);
      \draw[uaCharcoal!70,line width=0.9pt] (0,0) rectangle (4.5,3.6);
      % cells crossed by the facade
      \foreach \c in {(0.9,2.7),(1.8,2.7),(1.8,1.8),(2.7,1.8),(2.7,0.9),(3.6,0.9),(3.6,0)}
        \draw[uaOrange,line width=1.1pt] \c rectangle ++(0.9,0.9);
      % the facade itself
      \draw[uaCharcoal,line width=1.6pt] (1.35,3.6) -- (4.5,0.45);
      \node[text=uaCharcoal,rotate=-45] at (3.45,2.0) {\bfseries facade};
      \node[font=\footnotesize\bfseries,text=uaOrange] at (2.25,4.05)
        {Conservative cut cells};
      \node[text=uaGrey,align=center,text width=5.6cm] at (2.25,-0.55)
        {};
    \end{scope}
  \end{tikzpicture}
\end{frame}

% ---- PALM ------------------------------------------------------------
% \begin{frame}[t]{PALM}
%   \uacard[uaCharcoal]{\footnotesize
%     \uakeyt{Core.} Finite-difference Boussinesq. \\
%     \uakeyt{Geometry:} staircase. \\
%     \uakeyt{SGS:} Deardorff TKE.}
  % \vspace{-2pt}
  % \begin{columns}[T]
  %   \begin{column}{0.5\textwidth}
  %     {\bfseries\color{uaTeal}Advantages}\\[-3pt]
  %     {\color{uaTeal}\rule{\linewidth}{0.8pt}}
  %     \begin{itemize}\footnotesize\setlength\itemsep{2pt}
  %       \item high-order, accurate interior
  %       \item richest physics (radiation / chemistry / USM)
  %       \item scales to $\sim$\,50k cores; large domains
  %     \end{itemize}
  %   \end{column}
  %   \begin{column}{0.5\textwidth}
  %     {\bfseries\color{uaOrange}Drawbacks}\\[-3pt]
  %     {\color{uaOrange}\rule{\linewidth}{0.8pt}}
  %     \begin{itemize}\footnotesize\setlength\itemsep{2pt}
  %       \item staircase walls;
  %       \item upwind scheme is implicitly dissipative
  %       \item GPU only experimental; modules CPU-only
  %     \end{itemize}
  %   \end{column}
  % \end{columns}
% \end{frame}

% ---- PALM discretization (visual) -------------------------------------
\begin{frame}{PALM: staircase walls, upwind advection}
  \centering
  \begin{tikzpicture}[font=\scriptsize,>=latex]
    % ================= LEFT: staircase topography ======================
    % staircase fill (solid cells under the slanted surface)
    \fill[uaCharcoal!60] (0,0) rectangle (0.9,0.9);
    \fill[uaCharcoal!60] (0.9,0) rectangle (1.8,1.8);
    \fill[uaCharcoal!60] (1.8,0) rectangle (2.7,1.8);
    \fill[uaCharcoal!60] (2.7,0) rectangle (3.6,2.7);
    \fill[uaCharcoal!60] (3.6,0) rectangle (4.5,2.7);
    % MOST layer: first fluid cell above each step
    \fill[uaAmber!45] (0,0.9) rectangle (0.9,1.8);
    \fill[uaAmber!45] (0.9,1.8) rectangle (2.7,2.7);
    \fill[uaAmber!45] (2.7,2.7) rectangle (4.5,3.6);
    % grid
    \draw[uaCharcoal!45,line width=0.5pt,step=0.9] (0,0) grid (4.5,3.6);
    \draw[uaCharcoal!70,line width=0.9pt] (0,0) rectangle (4.5,3.6);
    % true surface
    \draw[uaOrange,line width=1.4pt,densely dashed] (0,1.0) -- (4.5,3.25);
    \node[text=uaOrange,rotate=26] at (1.6,2.05) {\bfseries true surface};
    \node[font=\footnotesize\bfseries,text=uaCharcoal] at (2.25,4.05)
      {Mask method: staircase walls};
    \node[text=uaGrey,align=center,text width=5.2cm] at (2.25,-0.65)
      {};

    % ================= RIGHT: 5th-order upwind stencil =================
    \begin{scope}[xshift=7.2cm]
      % wind direction
      \draw[->,uaTeal,line width=1.4pt] (0.4,3.3) -- (2.2,3.3)
        node[right,text=uaTeal]{\bfseries wind};
      % row of cells
      \draw[uaCharcoal!45,line width=0.6pt] (0,1.2) grid[step=0.9] (5.4,2.1);
      \draw[uaCharcoal!70,line width=0.9pt] (0,1.2) rectangle (5.4,2.1);
      % cell-center dots: 6-point stencil (4 upwind, 2 downwind)
      \foreach \x in {0.45,1.35,2.25,3.15}
        \fill[uaOrange] (\x,1.65) circle (2.6pt);
      \foreach \x in {4.05,4.95}
        \fill[uaOrange!40] (\x,1.65) circle (2.6pt);
      % the flux face
      \draw[uaTeal,line width=2.0pt] (3.6,1.08) -- (3.6,2.22);
      \node[text=uaTeal,anchor=north] at (3.6,1.0) {\bfseries flux face};
      % stencil brace
      \draw[uaOrange,line width=0.9pt] (0.45,2.45) -- (0.45,2.6) -- (4.95,2.6)
        -- (4.95,2.45);
      \node[text=uaOrange] at (2.7,2.85) {6-point stencil, upwind-biased};
      \node[font=\footnotesize\bfseries,text=uaCharcoal] at (2.7,4.05)
        {5th-order upwind advection};
      \node[text=uaGrey,align=center,text width=5.6cm] at (2.7,-0.65) {};
    \end{scope}
  \end{tikzpicture}
\end{frame}

% ---- LBM -------------------------------------------------------------
% \begin{frame}[t]{LBM (Evensen)}
%   \uacard[uaCharcoal]{\footnotesize
%     \uakeyt{Core.} Mesoscopic lattice-BGK. \\ 
%     \uakeyt{Geometry:} bounce-back. \\
%     \uakeyt{SGS:} Vreman. GPU-native.}
  % \vspace{-2pt}
  % \begin{columns}[T]
  %   \begin{column}{0.5\textwidth}
  %     {\bfseries\color{uaTeal}Advantages}\\[-3pt]
  %     {\color{uaTeal}\rule{\linewidth}{0.8pt}}
  %     \begin{itemize}\footnotesize\setlength\itemsep{2pt}
  %       \item GPU
  %       \item explicit local step
  %       \item density
  %     \end{itemize}
  %   \end{column}
  %   \begin{column}{0.5\textwidth}
  %     {\bfseries\color{uaOrange}Drawbacks}\\[-3pt]
  %     {\color{uaOrange}\rule{\linewidth}{0.8pt}}
  %     \begin{itemize}\footnotesize\setlength\itemsep{2pt}
  %       \item weakly compressible
  %       \item uniform lattice; bounce-back walls
  %       \item DA state choice: $f$ vs.\ macrovars
  %     \end{itemize}
  %   \end{column}
  % \end{columns}
% \end{frame}

% ---- LBM discretization (visual) --------------------------------------
\begin{frame}{LBM: collide \& stream on a uniform lattice}
  \centering
  \begin{tikzpicture}[font=\scriptsize,>=latex]
    % ================= LEFT: discrete velocity set (D2Q9) =============
    \draw[uaCharcoal!25,line width=0.5pt,step=1.1] (-1.1,-1.1) grid (1.1,1.1);
    \fill[uaCharcoal] (0,0) circle (3pt);
    \foreach \a in {0,90,180,270}
      \draw[->,uaTeal,line width=1.2pt] (0,0) -- (\a:1.0);
    \foreach \a in {45,135,225,315}
      \draw[->,uaTeal!60,line width=1.2pt] (0,0) -- (\a:1.35);
    \node[text=uaTeal] at (1.45,0.42) {$f_i$};
    \node[font=\footnotesize\bfseries,text=uaCharcoal] at (0,1.85)
      {Populations $f_i$};
    \node[text=uaGrey,align=center,text width=3.4cm] at (0,-1.95)
      {discrete velocity set\\ (D2Q9 sketch)};

    % ================= MIDDLE: collide & stream ========================
    \begin{scope}[xshift=4.6cm]
      \draw[uaCharcoal!25,line width=0.5pt,step=1.1] (-1.1,-1.1) grid (2.2,1.1);
      \foreach \x in {-1.1,0,1.1,2.2}
        \foreach \y in {-1.1,0,1.1}
          \fill[uaCharcoal!55] (\x,\y) circle (1.8pt);
      % stream arrows from node (0,0) to its neighbours
      \fill[uaCharcoal] (0,0) circle (3pt);
      \foreach \t in {(1.1,0),(0,1.1),(-1.1,0),(0,-1.1),(1.1,1.1),(-1.1,1.1),(1.1,-1.1),(-1.1,-1.1)}
        \draw[->,uaOrange,line width=1.1pt,shorten >=3pt,shorten <=4pt] (0,0) -- \t;
      \node[font=\footnotesize\bfseries,text=uaOrange] at (0.55,1.85)
        {Collide \& stream};
      \node[text=uaGrey,align=center,text width=4.2cm] at (0.55,-1.95)
        {local relaxation, then hop\\ to neighbours -- no global solve};
    \end{scope}

    % ================= RIGHT: bounce-back wall =========================
    \begin{scope}[xshift=10.1cm]
      % wall
      \fill[uaCharcoal!60] (-1.4,-1.5) rectangle (1.4,-1.1);
      \draw[uaCharcoal,line width=1.2pt] (-1.4,-1.1) -- (1.4,-1.1);
      \draw[uaCharcoal!25,line width=0.5pt,step=1.1] (-1.1,-1.1) grid (1.1,1.1);
      \fill[uaCharcoal] (0,0) circle (3pt);
      % incoming population toward wall, reflected back (offset for clarity)
      \draw[->,uaTeal,line width=1.2pt] (0.10,0.02) -- (-0.40,-0.98);
      \draw[->,uaOrange,line width=1.2pt,densely dashed] (-0.62,-0.98) -- (-0.12,0.02);
      \node[font=\footnotesize\bfseries,text=uaCharcoal] at (0,1.85)
        {Bounce-back walls};
      \node[text=uaGrey,align=center,text width=3.8cm] at (0,-2.35)
        {populations reflect\\ at solid nodes};
    \end{scope}
  \end{tikzpicture}
\end{frame}

% =====================================================================
%  PART 2 -- PYURBANAIR
% =====================================================================
\uasection[Part 2]{The \texttt{pyurbanair} framework}

\begin{frame}{Three Fortran solvers}
      uDALES, PALM and the LBM code are all \uakey{Fortran}
      \medskip
      \begin{itemize}\setlength\itemsep{8pt}
        \item \uakeyt{Compilers \& libraries} -- each needs its own
              build (MPI, netCDF\dots)
        \item \uakeyt{Config \& geometry} -- a different format for each
        \item \uakeyt{Output} -- different data types \& layouts
        \item \uakeyt{Warm start} -- handled differently by each
      \end{itemize}
  %   \begin{column}{0.44\textwidth}
  %     \uacard[uaOrange]{\footnotesize\centering
  %       \textbf{uDALES} $\cdot$ \textbf{PALM} $\cdot$ \textbf{LBM}\\[6pt]
  %       Cross-model experiments mean redoing all
  % of this, by hand, for every model.}
  %   \end{column}
\end{frame}

\begin{frame}{\texttt{pyurbanair}: one Python layer over all three}
  \begin{itemize}\setlength\itemsep{8pt}
    \item Automatic download
    \item Environment and library handling
    \item Unified config \& geometry
    \item Unified input/output
    \item Data assimilation
  \end{itemize}
\end{frame}

\begin{frame}{Architecture: one abstraction, many backends}
  \centering
  \vspace{-6pt}
  \begin{tikzpicture}[font=\scriptsize,>=latex,scale=0.94,transform shape,
      bk/.style={rounded corners=2pt,draw=uaTeal,line width=0.7pt,fill=white,
                 align=center,inner sep=3pt,text width=1.9cm,minimum height=0.7cm},
      top/.style={rounded corners=3pt,line width=0.9pt,align=center,inner sep=5pt},
      bar/.style={rounded corners=3pt,draw=uaOrange,line width=1pt,fill=uaOrange!10,
                  align=center,inner sep=5pt,text width=8.2cm,minimum height=0.5cm},
      da/.style={rounded corners=3pt,draw=uaTeal,line width=1pt,fill=uaTeal!16,
                 align=center,inner sep=5pt,text width=8.2cm,minimum height=0.5cm},
      ar/.style={->,line width=0.9pt,uaGrey}]
    % --- enclosing forward-model layer ---
    \draw[rounded corners=5pt,draw=uaTeal,line width=1.1pt,fill=uaTeal!5]
      (-0.5,0.65) rectangle (9.0,3.35);
    \node[font=\scriptsize\bfseries,text=uaTeal] at (4.25,3.07)
      {Forward-model layer};
    % backends
    \node[bk] (b1) at (0.75,2.35) {\textbf{pylbm}\\LBM};
    \node[bk] (b2) at (3.05,2.35) {\textbf{pyudales}\\uDALES};
    \node[bk] (b3) at (5.35,2.35) {\textbf{pypalm}\\PALM};
    \node[bk] (b4) at (7.65,2.35) {\textbf{neural}\\surrogate};
    % common interface bar
    \node[bar] (base) at (4.25,1.15)
      {\textbf{\color{uaOrange}BaseForwardModel} -- common interface};
    \draw[ar] (b1) -- (b1|-base.north);
    \draw[ar] (b2) -- (b2|-base.north);
    \draw[ar] (b3) -- (b3|-base.north);
    \draw[ar] (b4) -- (b4|-base.north);
    % --- common config + common setup ---
    \node[top,draw=uaAmber,fill=uaAmber!16] (cfg) at (1.85,4.20)
      {\textbf{\color{uaCharcoal}Common config}};
    \node[top,draw=uaTeal,fill=uaTeal!10] (setup) at (6.55,4.20)
      {\textbf{\color{uaCharcoal}Common setup:}\\ parameters \& initial states};
    \draw[ar] (cfg.south) -- (1.85,3.35);
    \draw[ar] (setup.south) -- (6.55,3.35);
    % --- data assimilation ---
    \node[da] (da) at (4.25,-0.35)
      {\textbf{\color{uaTeal}data-assimilation} (JAX) -- ESMDA $\cdot$
       localization $\cdot$ obs.\ operator};
    \draw[ar] (4.25,0.65) -- (da.north);
  \end{tikzpicture}
\end{frame}

% =====================================================================
%  PART 3 -- DATA ASSIMILATION: PARAMETER ESTIMATION
% =====================================================================
\uasection[Part 3]{Data assimilation: \\ parameter estimation}

% ---- pipeline / twin-experiment diagram ------------------------------
\begin{frame}{Setup}
  \centering
  \begin{tikzpicture}[font=\scriptsize,>=latex,scale=1.0,transform shape,
      every node/.append style={align=center},
      obox/.style={rounded corners=3pt,draw=uaOrange,line width=0.9pt,fill=uaOrange!6,
                   inner sep=4pt,text width=1.9cm,minimum height=0.85cm},
      tbox/.style={rounded corners=3pt,draw=uaTeal,line width=0.9pt,fill=uaTeal!6,
                   inner sep=4pt,text width=1.9cm,minimum height=0.85cm},
      ar/.style={->,line width=0.9pt}]
    % truth lane (top)
    \node[obox] (t1) at (0,1.5)    {truth inflow\\$\param(t)$};
    \node[obox] (t2) at (2.45,1.5) {truth\\solver};
    \node[obox] (t3) at (4.90,1.5) {truth\\state};
    \node[obox] (t4) at (7.35,1.5) {obs.\\operator};
    \node[obox,fill=uaOrange!16] (t5) at (9.80,1.5) {obs.\ $\obs$\\{\scriptsize($+$ noise)}};
    \draw[ar,uaOrange] (t1)--(t2);
    \draw[ar,uaOrange] (t2)--(t3);
    \draw[ar,uaOrange] (t3)--(t4);
    \draw[ar,uaOrange] (t4)--(t5);
    % assimilation lane (bottom)
    \node[tbox] (a1) at (0,-1.5)
      {prior inflow\\ensemble\\{\scriptsize($N_e=64$)}};
    \node[tbox,fill=uaTeal!10,text width=6.4cm,minimum height=1.1cm] (a2) at (5.2,-1.5)
      {\textbf{\color{uaTeal}ESMDA}\\[2pt]
       {\scriptsize forecast $\to$ predict obs.\ $\to$ update $\to$ repeat}};
    \node[tbox] (a3) at (9.80,-1.5) {posterior\\inflow};
    \draw[ar,uaTeal] (a1)--(a2);
    \draw[ar,uaTeal] (a2)--(a3);
    % observations feed INTO ES-MDA
    \draw[ar,uaCharcoal] (t5.south) --
       node[pos=0.62,right=1pt,font=\scriptsize]{observations} (a2.north east);
    % rollout loop
    \draw[ar,uaTeal] (a3.south) .. controls (9.8,-2.75) and (5.2,-2.75) ..
       node[pos=0.5,below,font=\scriptsize]{rollout: warm-start state, extrapolate prior} (a2.south);
  \end{tikzpicture}
\end{frame}

% ---- observation operator --------------------------------------------
\begin{frame}{The observation operator}
  \begin{columns}[T]
    \begin{column}{0.42\textwidth}
      \uakeyt{Spatial}: $N_{obs}$ sensors, \uakey{interpolated} from the grid
      \par\medskip
      \centering
      \begin{tikzpicture}[font=\tiny,scale=1.35,transform shape]
        \draw[uaCharcoal!22,line width=0.4pt] (0,0) grid (2,2);
        \foreach \x in {0,1,2}\foreach \y in {0,1,2}
          \fill[uaCharcoal!45] (\x,\y) circle (1.5pt);
        \coordinate (s) at (0.6,0.55);
        \draw[uaOrange!70,line width=0.6pt,densely dashed] (s)--(0,0);
        \draw[uaOrange!70,line width=0.6pt,densely dashed] (s)--(1,0);
        \draw[uaOrange!70,line width=0.6pt,densely dashed] (s)--(0,1);
        \draw[uaOrange!70,line width=0.6pt,densely dashed] (s)--(1,1);
        \fill[uaOrange] (s) circle (2.4pt);
        \node[text=uaOrange,anchor=south west,inner sep=1pt] at (s) {sensor};
        \node[text=uaGrey,anchor=north] at (1,-0.2) {grid nodes};
      \end{tikzpicture}
    \end{column}
    \begin{column}{0.58\textwidth}
      \centering
      \uakeyt{Temporal}: mean over \uakey{intervals}
      \par\medskip
      \begin{tikzpicture}[font=\tiny,>=latex,scale=1.4,transform shape]
        % bin edges
        \foreach \x in {0.5,1.0,...,4.0}
          \draw[uaCharcoal!25,line width=0.5pt,densely dashed] (\x,0.55) -- (\x,2.0);
        % axes
        \draw[->,uaCharcoal!70,line width=0.7pt] (0,0.55) -- (4.75,0.55);
        \node[text=uaGrey,anchor=north] at (0,0.5) {0};
        \node[text=uaGrey,anchor=north] at (4.5,0.5) {180 s};
        % 20 s bracket over first bin
        \draw[uaGrey,line width=0.6pt] (0,2.12) -- (0,2.22) -- (0.5,2.22) -- (0.5,2.12);
        \node[text=uaGrey,anchor=south] at (0.25,2.24) {10 s};
        % sensor signal (1 s cadence)
        \draw[uaTeal,line width=0.9pt] plot[smooth] coordinates
          {(0,1.30) (0.25,1.05) (0.5,0.85) (0.75,1.00) (1.0,1.35)
           (1.25,1.55) (1.5,1.45) (1.75,1.20) (2.0,1.10) (2.25,1.25)
           (2.5,1.50) (2.75,1.65) (3.0,1.55) (3.25,1.30) (3.5,1.05)
           (3.75,0.90) (4.0,1.00) (4.25,1.20) (4.5,1.35)};
        % per-bin means -> observations
        \foreach \xa/\m in {0.0/1.07, 0.5/1.07, 1.0/1.45, 1.5/1.25,
                            2.0/1.28, 2.5/1.57, 3.0/1.30, 3.5/0.98, 4.0/1.18}{
          \draw[uaOrange,line width=1.4pt] (\xa+0.05,\m) -- (\xa+0.45,\m);
          \fill[uaOrange] (\xa+0.25,\m) circle (1.6pt);}
        % legend
        \draw[uaTeal,line width=0.9pt] (0.1,0.15) -- (0.45,0.15);
        \node[text=uaGrey,anchor=west] at (0.5,0.15) {sensor signal};
        \draw[uaOrange,line width=1.4pt] (2.3,0.15) -- (2.65,0.15);
        \fill[uaOrange] (2.475,0.15) circle (1.6pt);
        \node[text=uaGrey,anchor=west] at (2.7,0.15) {bin mean = one obs.};
      \end{tikzpicture}
    \end{column}
  \end{columns}
\end{frame}


% ---- ESMDA core ------------------------------------------------------
\begin{frame}[t]{Ensemble Smoother with Multiple Data Assimilation (ESMDA)}
  \[
    \underbrace{\bm{m}_j^{a}}_{\color{uaGrey}\text{\scriptsize posterior}}
    =\underbrace{\bm{m}_j^{f}}_{\color{uaGrey}\text{\scriptsize prior}}
    +\overbrace{\mathbf{C}_{MD}\big(\mathbf{C}_{DD}+\alpha\,\mathbf{C}_D\big)^{-1}}
       ^{\color{uaTeal}\text{\scriptsize Kalman gain (from the ensemble)}}
    \big(\,\underbrace{\bm{d}+\sqrt{\alpha}\,\mathbf{C}_D^{1/2}\bm{z}_j}
              _{\color{uaOrange}\text{\scriptsize perturbed obs.}}
     -\underbrace{g(\bm{m}_j^{f})}_{\color{uaOrange}\text{\scriptsize predicted obs.}}\,\big)
  \]
  \vspace{-3pt}
  \begin{center}
    {\footnotesize $\displaystyle\sum_i \tfrac{1}{\alpha_i}=1$ -- the
    \uakeyt{inflation schedule} over $N_a$ steps.}
  \end{center}
\end{frame}

% ---- what ES-MDA updates ---------------------------------------------
\begin{frame}{What ESMDA actually updates}
  The inflow is a \uakey{trajectory}, carried on a few time knots per
  window:
  \[
    \bm{m}=\big[\,\underbrace{\alpha_0,\alpha_1,\dots}_{\text{inflow angle}},\;\;
    \underbrace{|U|_0,|U|_1,\dots}_{\text{wind speed}}\,\big].
  \]
  \begin{itemize}\setlength\itemsep{6pt}
    \item Smooth, time-correlated between knots
    \item Assimilate whole window at once
  \end{itemize}
\end{frame}


% ---- multi-window rollout --------------------------------------------
\begin{frame}{Windowed rollout: re-open uncertainty, keep the flow}
  \vspace{1pt}
  \centering
  \begin{tikzpicture}[font=\scriptsize,>=latex,yscale=0.92]
    % --- window boundaries, labels, time axis ---
    \foreach \x in {0,3,6,9}{\draw[uaGrey,densely dashed,line width=0.6pt] (\x,0.15) -- (\x,4.35);}
    \node[text=uaGrey] at (1.5,4.62) {window 1};
    \node[text=uaGrey] at (4.5,4.62) {window 2};
    \node[text=uaGrey] at (7.5,4.62) {window 3};
    \draw[->,uaGrey] (0,0.15) -- (9.5,0.15) node[right,text=black] {time};

    % ===== TOP: inflow parameter ensemble =====
    \node[text=uaAmber!85!black,align=center] at (-0.75,3.5) {inflow\\param.};
    % prior: wide ensemble -- tight at the seam, fanning open to the right
    \foreach \a in {0,3,6}{\fill[uaAmber!24]
        (\a,3.62) -- (\a+2.75,2.78) -- (\a+2.75,4.22) -- (\a,3.38) -- cycle;}
    % posterior: narrower (blue) band around the estimate
    \draw[uaTeal,opacity=0.32,line width=6.5pt,line cap=round]
      (0,3.5) .. controls (1.3,3.72) and (2.6,3.34) .. (3,3.5)
      .. controls (4.3,3.72) and (5.6,3.34) .. (6,3.5)
      .. controls (7.3,3.72) and (8.4,3.34) .. (9,3.52);
    % estimate / truth (black)
    \draw[black,line width=1pt]
      (0,3.5) .. controls (1.3,3.72) and (2.6,3.34) .. (3,3.5)
      .. controls (4.3,3.72) and (5.6,3.34) .. (6,3.5)
      .. controls (7.3,3.72) and (8.4,3.34) .. (9,3.52);
    % labels
    \node[text=uaAmber!85!black] at (2.5,2.48) {prior re-opens};
    \draw[->,uaAmber!85!black,line width=0.7pt] (2.15,2.62) -- (1.7,3.0);
    \node[text=uaTeal] at (5.5,2.72) {posterior};
    \draw[->,uaTeal,line width=0.7pt] (5.4,2.86) -- (5.15,3.42);

    % --- coupling: updated params re-simulate the state each update ---
    \node[text=uaGrey,align=center,font=\scriptsize\itshape] at (-0.75,2.45) {re-run\\model};
    \foreach \c in {1.5,4.5,7.5}{\draw[->,uaGrey,line width=0.6pt] (\c,2.66) -- (\c,2.12);}

    % ===== BOTTOM: turbulent state =====
    \node[text=uaTeal,align=center] at (-0.75,1.35) {turbulent\\state};
    % ensemble spread (re-simulated from the parameter ensemble)
    \draw[uaTeal,opacity=0.22,line width=8pt,line cap=round]
      (0,1.35) .. controls (0.5,1.75) and (1.0,0.85) .. (1.5,1.45)
      .. controls (2.0,2.0)  and (2.5,0.95) .. (3,1.55)
      .. controls (3.5,1.95) and (4.0,0.9)  .. (4.5,1.4)
      .. controls (5.0,1.85) and (5.5,0.95) .. (6,1.55)
      .. controls (6.5,1.95) and (7.0,0.85) .. (7.5,1.45)
      .. controls (8.0,1.9)  and (8.5,1.0)  .. (9,1.45);
    % mean state (thin); IC continuous across seams
    \draw[uaTeal,line width=1pt]
      (0,1.35) .. controls (0.5,1.75) and (1.0,0.85) .. (1.5,1.45)
      .. controls (2.0,2.0)  and (2.5,0.95) .. (3,1.55)
      .. controls (3.5,1.95) and (4.0,0.9)  .. (4.5,1.4)
      .. controls (5.0,1.85) and (5.5,0.95) .. (6,1.55)
      .. controls (6.5,1.95) and (7.0,0.85) .. (7.5,1.45)
      .. controls (8.0,1.9)  and (8.5,1.0)  .. (9,1.45);
    \node[text=uaTeal] at (4.5,0.5) {re-simulated from the parameters};
  \end{tikzpicture}
  \par\smallskip
  {\footnotesize\color{uaGrey}Every update \uakeyt{re-simulates} the state
  from the new parameters -- only the IC carries over.}
\end{frame}

% ---- AR(2) dynamic parameters ----------------------------------------
\begin{frame}{Dynamic parameters: AR(2) trajectories}
  \par
  \vspace{20pt}
  \centering
  \begin{tikzpicture}[font=\scriptsize,>=latex,xscale=1.5]
    % external envelope band (new window)
    \fill[uaAmber!16] (0,0.9) rectangle (4.6,2.1);
    \draw[uaGrey,densely dashed,line width=0.7pt] (-1.5,1.5) -- (4.6,1.5);
    \node[text=uaGrey,anchor=west,align=left] at (4.62,1.5)
      {$x_{\text{ext}}$};
    \draw[<->,uaAmber!80!black,line width=0.8pt] (4.38,0.9) -- (4.38,2.1);
    \node[text=uaAmber!80!black,anchor=east] at (4.34,1.78)
      {$\pm\Sigma_{\text{ext}}$};
    % previous-window posterior (tight bundle, orange)
    \foreach \o in {-0.07,0,0.07}
      \draw[uaOrange,line width=0.8pt] plot[smooth] coordinates
        {(-1.5,1.85+\o) (-1.0,2.0+\o) (-0.5,2.25+\o) (0,2.6)};
    \node[text=uaOrange,align=center] at (-0.78,2.85)
      {posterior,\\ window $w$};
    % window seam
    \draw[uaCharcoal!60,densely dashed,line width=0.8pt] (0,0.55) -- (0,3.35);
    % new-window prior ensemble (teal), pinned in value & slope at the seam
    \draw[uaTeal,line width=0.8pt] plot[smooth] coordinates
      {(0,2.6) (0.5,2.95) (1,3.05) (1.5,2.95) (2,2.7) (2.5,2.45) (3,2.25) (3.5,2.15) (4.2,2.1)};
    \draw[uaTeal,line width=0.8pt] plot[smooth] coordinates
      {(0,2.6) (0.5,2.92) (1,3.0) (1.5,2.8) (2,2.5) (2.5,2.2) (3,2.0) (3.5,1.9) (4.2,1.95)};
    \draw[uaTeal,line width=0.8pt] plot[smooth] coordinates
      {(0,2.6) (0.5,2.93) (1,2.95) (1.5,2.7) (2,2.35) (2.5,2.05) (3,1.85) (3.5,1.7) (4.2,1.65)};
    \draw[uaTeal,line width=0.8pt] plot[smooth] coordinates
      {(0,2.6) (0.5,2.9) (1,2.9) (1.5,2.6) (2,2.2) (2.5,1.9) (3,1.7) (3.5,1.6) (4.2,1.5)};
    \draw[uaTeal,line width=0.8pt] plot[smooth] coordinates
      {(0,2.6) (0.5,2.88) (1,2.85) (1.5,2.5) (2,2.1) (2.5,1.7) (3,1.4) (3.5,1.2) (4.2,1.1)};
    \draw[uaTeal,line width=0.8pt] plot[smooth] coordinates
      {(0,2.6) (0.5,2.9) (1,2.8) (1.5,2.4) (2,1.95) (2.5,1.55) (3,1.3) (3.5,1.05) (4.2,0.95)};
    \fill[uaCharcoal] (0,2.6) circle (1.8pt);
    % annotations
    \node[text=uaCharcoal,align=center,anchor=west] at (0.65,3.5)
      {\bfseries pinned at the seam: value \emph{and} slope};
    \draw[->,uaCharcoal,line width=0.7pt] (0.62,3.45) -- (0.06,2.72);
    \node[text=uaTeal,align=center] at (2.9,3.2)
      {prior ensemble, window $w{+}1$};
    \node[text=uaGrey,align=center] at (2.3,0.72)
      {relaxes to $x_{\text{ext}}$; spread re-opens toward $\Sigma_{\text{ext}}$};
    % t axis
    \draw[->,uaCharcoal!70,line width=0.7pt] (-1.5,0.55) -- (4.7,0.55);
    \node[text=uaGrey,anchor=north] at (4.6,0.5) {$t$};
  \end{tikzpicture}
\end{frame}

% =====================================================================
%  GROUND TRUTH & SENSORS  (right before results)
% =====================================================================

\begin{frame}{Sensor observations}
  \centering
  \animategraphics[autoplay,loop,controls,width=0.9\linewidth]{9}%
    {figures/anim_sensors_frames/g_}{0}{67}
  \par\smallskip
  {\footnotesize\color{uaGrey}Sensors sample $u,v,w$; each is \uakey{averaged
  over 10\,s} into one observation. \textit{(plays in Adobe Acrobat)}}
\end{frame}

% =====================================================================
%  PART 3 -- PARAMETER-ESTIMATION RESULTS  (Block A)
% =====================================================================

% ===== Block A: parameter-only, model x resolution =====
% ---- parameter trajectories at highest resolution ($\Delta x = 1$ m) ----
\begin{frame}{uDALES ($\Delta x = 1$ m): parameter trajectories}
  \centering
  \includegraphics[width=0.72\linewidth]{figures/A_param_only/results/param_traj_udales.png}
  \par\smallskip
  {\footnotesize\color{uaGrey}Inflow angle \& speed: \uakeyt{posterior} (blue) vs.\
  \uakey{truth} (black dashed); prior in orange.}
\end{frame}

\begin{frame}{uDALES ($\Delta x = 1$ m): reconstructed final state}
  \centering
  \includegraphics[width=\linewidth]{figures/A_param_only/results/final_state_udales.png}
  \par\smallskip
  {\footnotesize\color{uaGrey}$|U|$ at the final time: \uakey{truth},
  \uakeyt{posterior mean}, and ensemble std; dots = assimilation sensors.}
\end{frame}

\begin{frame}{PALM ($\Delta x = 1$ m): parameter trajectories}
  \centering
  \includegraphics[width=0.72\linewidth]{figures/A_param_only/results/param_traj_palm.png}
  \par\smallskip
  {\footnotesize\color{uaGrey}Inflow angle \& speed: \uakeyt{posterior} (blue) vs.\
  \uakey{truth} (black dashed); prior in orange.}
\end{frame}


\begin{frame}{PALM ($\Delta x = 1$ m): reconstructed final state}
  \centering
  \includegraphics[width=\linewidth]{figures/A_param_only/results/final_state_palm.png}
  \par\smallskip
  {\footnotesize\color{uaGrey}$|U|$ at the final time: \uakey{truth},
  \uakeyt{posterior mean}, and ensemble std; dots = assimilation sensors.}
\end{frame}


\begin{frame}{LBM ($\Delta x = 1$ m): parameter trajectories}
  \centering
  \includegraphics[width=0.72\linewidth]{figures/A_param_only/results/param_traj_lbm.png}
  \par\smallskip
  {\footnotesize\color{uaGrey}Inflow angle \& speed: \uakeyt{posterior} (blue) vs.\
  \uakey{truth} (black dashed); prior in orange.}
\end{frame}

\begin{frame}{LBM ($\Delta x = 1$ m): reconstructed final state}
  \centering
  \includegraphics[width=\linewidth]{figures/A_param_only/results/final_state_lbm.png}
  \par\smallskip
  {\footnotesize\color{uaGrey}$|U|$ at the final time: \uakey{truth},
  \uakeyt{posterior mean}, and ensemble std; dots = assimilation sensors.}
\end{frame}

\begin{frame}{Accuracy at a glance: parameters \& state}
  \centering\small
  \renewcommand{\arraystretch}{1.15}
  \begin{tabular}{llcccc}
    \toprule
    \rowcolor{uaTeal!12}
    \textbf{Model} & \textbf{$\Delta x$} & \textbf{$\alpha$ RMSE} &
      \textbf{$|U|$ RMSE} & \textbf{field RMSE} & \textbf{val.\ RMSE}\\
    \rowcolor{uaTeal!12}
     & [m] & [deg] & [m/s] & [m/s] & [m/s]\\
    \midrule
    LBM    & 2   & 27.0 & 2.23 & 2.34 & 1.54\\
    LBM    & 1.3 & 11.2 & 0.61 & 0.78 & 0.35\\
    LBM    & 1   & 14.2 & 0.90 & 1.00 & 0.78\\
    PALM   & 2   & 11.8 & 0.68 & 0.96 & 0.84\\
    PALM   & 1.3 & 15.5 & 0.62 & 0.78 & 0.63\\
    PALM   & 1   & 11.1 & 0.35 & 0.76 & 1.00\\
    uDALES & 2   & 18.5 & 1.18 & 1.01 & 0.71\\
    uDALES & 1.3 & 7.49 & 0.45 & 0.37 & 0.31\\
    \rowcolor{uaAmber!18}
    \textbf{uDALES} & \textbf{1} & \textbf{3.37} & \textbf{0.30} &
      \textbf{0.26} & \textbf{0.24}\\
    \bottomrule
  \end{tabular}
\end{frame}

% =====================================================================
%  PART 4 -- DATA ASSIMILATION: STATE & PARAMETER ESTIMATION
% =====================================================================
\uasection[Part 4]{\\ Data assimilation: \\
state \& parameter \\estimation}

\begin{frame}{Estimating the state, not just the inflow}
      \smallskip

      ESMDA now acts on the \uakeyt{augmented state}
      $\bm{w}=[\,\bm{x},\,\bm{\theta}\,]$.
      \medskip
      \begin{itemize}\setlength\itemsep{7pt}
        \item The state is \uakey{high-dimensional} ($10^5$--$10^7$)
        \item Two routes: \uakeyt{localize} the update, or \uakeyt{reduce} the
              state dimension
        \item Focus on initial condition in each window
      \end{itemize}
\end{frame}

\begin{frame}[t]{Localizing the initial-condition update}
  \smallskip
  \begin{columns}[T]
    \begin{column}{0.5\textwidth}
      {\bfseries\color{uaOrange}Distance-based}\\[-3pt]{\color{uaOrange}\rule{\linewidth}{0.8pt}}
      \begin{itemize}\footnotesize\setlength\itemsep{4pt} 
        \item keep obs within a physical distance $d_t$
        \item must be \uakey{predefined}
      \end{itemize}
    \end{column}
    \begin{column}{0.5\textwidth}
      {\bfseries\color{uaTeal}Correlation-based (adaptive)}\\[-3pt]{\color{uaTeal}\rule{\linewidth}{0.8pt}}
      \begin{itemize}\footnotesize\setlength\itemsep{4pt}
        \item keep obs with ensemble correlation $|\rho|>\rho_t$
        \item \uakey{no coordinates}
      \end{itemize}
    \end{column}
  \end{columns}
\end{frame}

\begin{frame}{State reduction with an SVD subspace}
  \begin{columns}[T]
    \begin{column}{0.55\textwidth}
      Alternative to localization. Two variants:
      \begin{itemize}\setlength\itemsep{8pt}
        \item \uakeyt{reduced IC} -- estimate the window's initial condition
        \item \uakeyt{reduced IC + state} -- plus a \uakey{final full-state
              Kalman update}
      \end{itemize}
      \medskip
    \end{column}
    \begin{column}{0.45\textwidth}
      \centering
      \begin{tikzpicture}[font=\scriptsize,>=latex,scale=0.92,transform shape,
          big/.style={rounded corners=3pt,draw=uaTeal,line width=1pt,fill=uaTeal!8,
                      align=center,inner sep=4pt,text width=3.7cm,minimum height=0.9cm},
          red/.style={rounded corners=3pt,draw=uaOrange,line width=1.2pt,fill=uaOrange!12,
                      align=center,inner sep=4pt,text width=2.3cm,minimum height=0.95cm},
          dim/.style={text=uaGrey,font=\tiny,align=center}]
        % --- nodes (wide -> narrow -> wide) ---
        \node[big] (fs)  at (0,4.25) {full state $\bm{x}$\\{\tiny $u,v,w$ on the grid}};
        \node[dim] at (3.05,4.25) {$10^5$\\$-10^7$};
        \node[red] (z)   at (0,2.15) {\textbf{\color{uaOrange}ES-MDA in $\bm{z}$}\\{\tiny reduced coeffs}};
        \node[dim] at (-2.6,2.15) {dim\\$\ll$};
        \node[big] (fs2) at (0,0.05) {updated state $\bm{x}$};
        % --- encode / decode arrows ---
        \draw[->,line width=1pt,uaTeal] (fs.south) -- (z.north)
          node[midway,right=1pt,text=uaTeal] {$\bm{z}=\bm{\Phi}^{\!\top}\bm{x}$}
          node[midway,left=1pt,dim] {encode};
        \draw[->,line width=1pt,uaTeal] (z.south) -- (fs2.north)
          node[midway,right=1pt,text=uaTeal] {$\bm{x}=\bm{\Phi}\bm{z}$}
          node[midway,left=1pt,dim] {decode};
        % --- ES-MDA iteration self-loop ---
        \draw[->,uaOrange,line width=0.9pt]
          (z.east) .. controls +(1.05,0.6) and +(1.05,-0.6) .. (z.east);
        \node[text=uaOrange,font=\tiny,align=center] at (2.65,2.15) {$N_a$\\iter.};
      \end{tikzpicture}
      \par\smallskip
      {\tiny\color{uaGrey}$\bm{\Phi}$: SVD/POD basis of the ensemble}
    \end{column}
  \end{columns}
\end{frame}

% ===== state-estimation results (Block B) =====
% ===== Block B: does state estimation help? (uDALES) =====
% \begin{frame}{Does estimating the state help? (uDALES)}
%   \centering
%   \includegraphics[width=0.9\linewidth]{figures/B_state_estimation/B2_state_err_vs_time_methods.pdf}
%   \par\smallskip
%   {\footnotesize\color{uaGrey}\uakey{param-only} (black) stays lowest; state
%   methods \uakey{spike at each window start} (30\,s windows).}
% \end{frame}

\begin{frame}{Validation sensors by method (uDALES)}
  \centering
  \includegraphics[width=0.74\linewidth]{figures/B_state_estimation/B4_valsensors_methods.pdf}
  \par\smallskip
\end{frame}

\begin{frame}{Method comparison (uDALES, 0--300\,s)}
  \centering\small
  \renewcommand{\arraystretch}{1.25}
  \begin{tabular}{lccccc}
    \toprule
    \rowcolor{uaTeal!12}
    \textbf{Method} & \textbf{$\alpha$ RMSE} & \textbf{$|U|$ RMSE} &
      \textbf{field RMSE} & \textbf{val.\ RMSE} & \textbf{$N_e$}\\
    \rowcolor{uaTeal!12}
     & [deg] & [m/s] & [m/s] & [m/s] & \\
    \midrule
    \rowcolor{uaAmber!18}
    \textbf{param-only} & \textbf{1.02} & \textbf{0.25} & \textbf{0.11} &
      \textbf{0.08} & 64\\
    IC+param (corr)            & 6.34 & 0.75 & 0.52 & 0.35 & 128\\
    IC+param (dist)            & 6.35 & 0.75 & 0.51 & 0.41 & 128\\
    IC+param (reduction)       & 6.97 & 0.79 & 0.54 & 0.41 & 128\\
    IC+state+param (reduction) & 5.00 & 0.43 & 0.34 & 0.47 & 128\\
    \bottomrule
  \end{tabular}
\end{frame}

% \begin{frame}{Cost vs.\ accuracy}
%   \centering
%   \includegraphics[width=0.62\linewidth]{figures/summary/S1_cost_vs_accuracy.pdf}
%   \par\smallskip
%   {\footnotesize\color{uaGrey}\uakeyt{uDALES} gives the best accuracy per
%   wall-clock; LBM needs fine grids to compete.}
% \end{frame}

% =====================================================================
%  PART 5 -- NEURAL SURROGATES
% =====================================================================
\uasection[Part 5]{Neural surrogates}

\begin{frame}{Why a surrogate?}
  \begin{columns}[T]
    \begin{column}{0.56\textwidth}
      \begin{itemize}\setlength\itemsep{6pt}
        \item LES forward models are very expensive to run!
        \item We need to run them many times
        \item Train a surrogate \uakeyt{offline}, then use it \uakeyt{online}
      \end{itemize}
    \end{column}
    \begin{column}{0.44\textwidth}
      \uacard[uaOrange]{\footnotesize\centering
        \textbf{\color{uaOrange}ESMDA inner loop}\\[5pt]
        forecast $N_e$ members\\[5pt]
        $\downarrow$\\ predict observations\\ $\downarrow$\\
        update parameters\\ $\downarrow$\\ repeat\\[6pt]
        \uakeyt{Surrogate replaces the forecast}}
    \end{column}
  \end{columns}
\end{frame}

% ---- 2. general setup -------------------------------------------------
\begin{frame}{General setup: learn the one-step transition}
  \smallskip
  \par\centering
  \begin{tikzpicture}[font=\scriptsize,>=latex,scale=0.96,transform shape,
      inp/.style={rounded corners=3pt,draw=uaTeal,line width=0.9pt,fill=uaTeal!6,
                  align=center,inner sep=4pt,text width=2.5cm,minimum height=0.62cm},
      net/.style={rounded corners=4pt,draw=uaOrange,line width=1.1pt,fill=uaOrange!10,
                  align=center,inner sep=6pt,text width=3.1cm,minimum height=2.1cm},
      outp/.style={rounded corners=3pt,draw=uaTeal,line width=1pt,fill=uaTeal!16,
                  align=center,inner sep=4pt,text width=2.2cm,minimum height=0.8cm}]
    \node[inp] (st) at (0,1.45) {state $\state_t$\\{\tiny $u,v,w$ on the grid}};
    \node[inp] (pr) at (0,0.55) {inflow $\param_t$\\{\tiny angle, speed}};
    \node[inp] (ge) at (0,-0.45) {geometry mask\\{\tiny fluid/obstacle}};
    \node[net] (nn) at (4.5,0.5)
      {\textbf{\color{uaOrange}3D ConvNeXt U-Net}};
    \node[outp] (o) at (8.4,0.5) {state $\state_{t+1}$};
    \draw[->,line width=0.9pt,uaTeal] (st.east) -- (nn.west|-st.east);
    \draw[->,line width=0.9pt,uaTeal] (pr.east) -- (nn.west|-pr.east);
    \draw[->,line width=0.9pt,uaTeal] (ge.east) -- (nn.west|-ge.east);
    \draw[->,line width=1.1pt,uaOrange] (nn.east) -- (o.west);
  \end{tikzpicture}
\end{frame}

% ---- 3. architecture ---------------------------------------------------
% \begin{frame}{Architecture: 3D ConvNeXt-UNet, 2.38\,M parameters}
%   \begin{columns}[T]
%     \begin{column}{0.46\textwidth}
%       \centering
%       \begin{tikzpicture}[font=\scriptsize,
%           blk/.style={rounded corners=2pt,draw=uaTeal,line width=0.8pt,fill=uaTeal!6,
%                       align=center,inner sep=3pt,text width=3.3cm,minimum height=0.5cm}]
%         \node[font=\footnotesize\bfseries,text=uaCharcoal] at (0,4.45)
%           {One ConvNeXt block};
%         \node[blk] (b1) at (0,3.8) {depthwise conv $k{=}7$\\{\tiny separable: 3 axis-aligned passes}};
%         \node[blk] (b2) at (0,2.9)  {GroupNorm(8)};
%         \node[blk,draw=uaOrange,fill=uaOrange!10] (b3) at (0,2.1)
%           {\textbf{\color{uaOrange}FiLM}: scale \& shift\\{\tiny from the inflow $\param$}};
%         \node[blk] (b4) at (0,1.2) {$1{\times}1$ expand $\to$ GELU $\to$ project};
%         \foreach \a/\b in {b1/b2,b2/b3,b3/b4}
%           \draw[->,line width=0.8pt,uaGrey] (\a) -- (\b);
%         % residual
%         \draw[->,line width=0.8pt,uaGrey,rounded corners=3pt]
%           (0,4.18) -| (2.1,0.55) -- (0.4,0.55);
%         \node[circle,draw=uaGrey,inner sep=1pt,fill=white] at (0.25,0.55) {$+$};
%         \node[text=uaGrey,font=\tiny] at (2.55,2.35) {residual};
%       \end{tikzpicture}
%     \end{column}
%     \begin{column}{0.54\textwidth}
%       U-Net encoder--decoder with skips; stages of ConvNeXt blocks
%       {\footnotesize(base 32, depths $[3,3,3,3]$)}.
%       \medskip

%       \uakeyt{Physics-aware wrapping:}
%       \begin{itemize}\footnotesize\setlength\itemsep{4pt}
%         \item \uakey{Fluid masking} of inputs, outputs and loss --
%               uDALES dumps carry junk inside obstacles; solid cells
%               stay exactly zero in rollouts
%         \item \uakey{Z-score normalization} over fluid cells, stored as
%               model buffers -- ships with the weights
%         \item \uakey{Residual prediction}: the head outputs the state
%               \emph{increment} -- ``no change'' is the zero solution;
%               markedly more stable rollouts
%       \end{itemize}
%     \end{column}
%   \end{columns}
% \end{frame}

% ---- 3b. full architecture diagram --------------------------------------
\begin{frame}{Architecture: ConvNeXt U-Net}
  \centering
  \includegraphics[width=\linewidth]{figures/E_surrogate/drift_term_unet.pdf}
\end{frame}

% ---- 4. training --------------------------------------------------------
\begin{frame}{Training: pushforward unrolling + horizon curriculum}
  \par\centering
  \begin{tikzpicture}[font=\scriptsize,>=latex,
      s/.style={rounded corners=3pt,draw=uaCharcoal!60,line width=0.8pt,fill=white,
                align=center,inner sep=3pt,minimum height=0.6cm,text width=1.0cm},
      sp/.style={rounded corners=3pt,draw=uaGrey,line width=0.8pt,fill=uaCharcoal!6,
                 align=center,inner sep=3pt,minimum height=0.6cm,text width=1.0cm},
      sg/.style={rounded corners=3pt,draw=uaOrange,line width=1pt,fill=uaOrange!10,
                 align=center,inner sep=3pt,minimum height=0.6cm,text width=1.0cm},
      tr/.style={rounded corners=3pt,draw=uaTeal,line width=1pt,fill=uaTeal!10,
                 align=center,inner sep=3pt,minimum height=0.6cm,text width=1.0cm}]
    \node[s]  (x0) at (0,0)     {$\state_t$\\{\tiny data}};
    \node[sp] (x1) at (1.85,0)  {$\hat\state_{t+1}$};
    \node[sp] (x2) at (3.70,0)  {$\hat\state_{t+2}$};
    \node[sp] (x3) at (5.55,0)  {$\hat\state_{t+3}$};
    \node[sg] (x4) at (7.40,0)  {$\hat\state_{t+4}$};
    \node[sg] (x5) at (9.25,0)  {$\hat\state_{t+5}$};
    \node[tr] (xt) at (11.45,0) {$\state_{t+5}$\\{\tiny data}};
    \draw[->,line width=0.9pt,uaGrey]   (x0) -- (x1);
    \draw[->,line width=0.9pt,uaGrey]   (x1) -- (x2);
    \draw[->,line width=0.9pt,uaGrey]   (x2) -- (x3);
    \draw[->,line width=0.9pt,uaOrange] (x3) -- (x4);
    \draw[->,line width=0.9pt,uaOrange] (x4) -- (x5);
    \draw[<->,line width=0.9pt,uaTeal]  (x5) -- node[above]{loss} (xt);
    \node[text=uaGrey,font=\scriptsize] at (2.8,-0.85)
      {predictions, no gradients};
    \node[text=uaOrange,font=\scriptsize\bfseries] at (8.3,-0.85)
      {gradients through the last 2 steps};
  \end{tikzpicture}
  \par\medskip
  \raggedright
  \begin{itemize}\setlength\itemsep{4pt}
    \item \uakey{Rollout} predictions
    \item Only compute gradient of last 2 steps
    \item Ramp up rollout steps during training
  \end{itemize}
\end{frame}

% ---- 5. computational efficiency ----------------------------------------
\begin{frame}{Computational efficiency}
  \begin{itemize}\setlength\itemsep{8pt}
    \item Separable depthwise convs
    \item Compile model
    \item bf16 autocast
    \item Async dataloading
    \item Geometry de-duplication
  \end{itemize}
\end{frame}

% ---- 6. inside pyurbanair -------------------------------------------------
\begin{frame}{Inside \texttt{pyurbanair}: a drop-in forward model}
      \smallskip
      \begin{itemize}\setlength\itemsep{4pt}
        \item Trained neural network is wrapped in the same forward model setup
        \item Cold start: \uakeyt{uDALES spins up} the flow
        \item Vectorized over \uakeyt{ESMDA ensemble members}
      \end{itemize}
\end{frame}

\begin{frame}{Neural surrogate: parameter trajectories}
  \centering
  \includegraphics[width=0.72\linewidth]{figures/E_surrogate/param_traj_surrogate.png}
\end{frame}

\begin{frame}{Surrogate vs.\ truth: reconstructed state (0--600\,s)}
  \centering
  \animategraphics[autoplay,loop,controls,width=0.78\linewidth]{10}%
    {figures/anim_surrogate_truth/g_}{0}{59}
\end{frame}

% =====================================================================
%  NEXT STEPS  (animated state-estimation videos -- play in Adobe Acrobat)
% =====================================================================
\uasection[Part 7]{Next steps}

\begin{frame}{Where we go from here}
  \begin{itemize}\setlength\itemsep{14pt}
    \item \uakeyt{Unify model settings further}
    \item \uakeyt{More tests on hyperparameters} -- ensemble size,
          localization, window length, observation intervals
    \item \uakeyt{State-estimation methods} -- improve state localization and reduction
    \item \uakeyt{Improve the neural surrogate} -- accuracy and stability over
          longer rollouts
    \item \uakeyt{Larger cases} -- scale up to bigger domains and real urban
          geometries
  \end{itemize}
\end{frame}

\begin{frame}{Geometry from Barcelona: uDALES, PALM, LBM}
  \centering
  \animategraphics[autoplay,loop,controls,width=\linewidth]{10}%
    {figures/anim_state3/g_}{0}{89}
  \par\smallskip
  {\footnotesize\color{uaGrey} Estimated velocity magnitude $|U|$ over the
  building array -- the three solvers side by side.
  \textit{(animation plays in Adobe Acrobat)}}
\end{frame}

% =====================================================================
%  SUMMARY
% =====================================================================
\begin{frame}{Conclusion}
  \begin{itemize}\setlength\itemsep{10pt}
    \item Not easy to simply use a different model!
    \item State estimation is difficult!
    \item Neural surrogates show potential
  \end{itemize}
\end{frame}

% =====================================================================
%  THANKS
% =====================================================================
\begin{frame}[plain]
  \begin{tikzpicture}[remember picture,overlay]
    \fill[uaTeal](current page.south west)rectangle(current page.north east);
    \fill[white,opacity=0.10]($(current page.north east)$) arc (0:-180:5.2cm and 5.2cm);
    \node[anchor=north west,inner sep=0pt] at ($(current page.north west)+(\uaMargin,-0.4cm)$)
      {\includegraphics[height=1.0cm]{\uaLogoWhite}};
    \node[anchor=west,text width=11cm,inner sep=0pt] at ($(current page.west)+(\uaMargin,0.7cm)$)
      {\uaTitleFont\bfseries\color{white}\Huge Thank you!\par
       \vspace{0.3cm}\color{white}\normalsize Questions \& discussion welcome.\par};
    \node[anchor=west,text width=11cm,inner sep=0pt] at ($(current page.west)+(\uaMargin,-1.7cm)$)
      {\color{uaLightBlue}\footnotesize
       Nikolaj T.\ M\"ucke \;$\cdot$\; ntmucke@tudelft.nl \;$\cdot$\;
       github.com/nmucke/pyurbanair\par};
  \end{tikzpicture}
\end{frame}

\end{document}
